% Physical pages: 436--441 (printed pages 413--418).
% Figures: none.
% Uncertainties: none.

\subsection{Path-Integral Formulation of Quantum Electrodynamics}

The path-integral approach to quantum field theory really comes into its own
when applied to gauge theories of massless spin one particles, such as quantum
electrodynamics. The derivation of the Feynman rules for quantum
electrodynamics in the previous chapter involved a fair amount of hand-waving,
in arguing that the terms in the photon propagator $\Delta^{\mu\nu}(q)$
proportional to $q^\mu$ or $q^\nu$ could be dropped, and that the purely
time-like terms would just cancel the Coulomb term in the Hamiltonian, so that
the effective photon propagator could be taken as $\eta^{\mu\nu}/q^2$. To give
a real justification of this result using the methods of Chapter 8 would
involve us in a complicated analysis of Feynman diagrams. But as we shall now
see, the path-integral approach yields the desired form of the photon
propagator, without ever having to think about the details of Feynman diagrams.

In Chapter 8 we found that in Coulomb gauge, the Hamiltonian for the interaction
of photons with charged particles takes the form
\begin{equation}
H[\mathbf A,\boldsymbol\Pi_\perp,\ldots]
=H_M+\int d^3x\left[\frac12\boldsymbol\Pi_\perp^2
 +\frac12(\boldsymbol\nabla\times\mathbf A)^2-\mathbf A\cdot\mathbf J\right]
 +V_{\mathrm{Coul}}.
\tag{9.6.1}\label{eq:9.6.1}
\end{equation}
Here $\mathbf A$ is the vector potential, subject to the Coulomb gauge condition
\begin{equation}
\boldsymbol\nabla\cdot\mathbf A=0,
\tag{9.6.2}\label{eq:9.6.2}
\end{equation}
while $\boldsymbol\Pi_\perp$ is the solenoidal part of its canonical conjugate,
satisfying the same constraint
\begin{equation}
\boldsymbol\nabla\cdot\boldsymbol\Pi_\perp=0.
\tag{9.6.3}\label{eq:9.6.3}
\end{equation}
Also, $H_M$ is the matter Hamiltonian and $V_{\mathrm{Coul}}$ is the Coulomb
energy
\begin{equation}
V_{\mathrm{Coul}}(t)=\frac12\int d^3x\,d^3y\,
 \frac{J^0(\mathbf{x},t)J^0(\mathbf{y},t)}{4\pi|\mathbf{x}-\mathbf{y}|}.
\tag{9.6.4}\label{eq:9.6.4}
\end{equation}

Just as for any other Hamiltonian system, we can calculate vacuum expectation
values of time-ordered products as path integrals\footnote{Note that
$\boldsymbol\pi(x)$ is the interpolating c-number field for the quantum
operator $\boldsymbol\Pi_\perp$, whose commutation relations with each other
and with $\mathbf A$ are the same as those of $\boldsymbol\Pi$, but which unlike
$\boldsymbol\Pi$ commutes with all canonical matter variables.}
\begin{align}
\left\langle T\{\mathcal O_A\mathcal O_B\cdots\}\right\rangle_{\Omega}
={}&\int\left[\prod_{x,i}da_i(x)\prod_{x,i}d\pi_i(x)
 \prod_{x,\ell}d\psi_\ell(x)\right]
\notag\\
&\cdot\mathcal O_A\mathcal O_B\cdots
\notag\\
&\cdot\exp\Biggl\{i\int d^4x\Biggl[
 \boldsymbol\pi\cdot\dot{\mathbf a}-\frac12\boldsymbol\pi^2
 -\frac12(\boldsymbol\nabla\times\mathbf a)^2
\notag\\
&\qquad+\mathbf a\cdot\mathbf j+\mathcal L_M\Biggr]
 -i\int dt\,V_{\mathrm{Coul}}\Biggr\}
\notag\\
&\cdot\left[\prod_x\delta(\boldsymbol\nabla\cdot\mathbf a(x))\right]
 \left[\prod_x\delta(\boldsymbol\nabla\cdot\boldsymbol\pi(x))\right],
\tag{9.6.5}\label{eq:9.6.5}
\end{align}
where $\psi_\ell(x)$ are generic matter fields. In writing Eq. \eqref{eq:9.6.5} in
terms of a matter Lagrangian density, we are assuming that $H_M$ is local and
either linear in the matter $\pi$s (as in spinor electrodynamics) or quadratic
with field-independent coefficients (as in scalar electrodynamics). We have
inserted delta functions\footnote{This is not strictly accurate. If we take
the canonical variables to be, say, $a_1,a_2$ and $\pi_1,\pi_2$, with $a_3$
and $\pi_3$ regarded as functionals of these variables given by Eqs. \eqref{eq:9.6.2}
and \eqref{eq:9.6.3}, then we should insert the delta functions
\[
\prod_x\delta\!\left(a_3(x)+\partial_3^{-1}
 [\partial_1a_1(x)+\partial_2a_2(x)]\right)
\delta\!\left(\pi_3(x)+\partial_3^{-1}
 [\partial_1\pi_1(x)+\partial_2\pi_2(x)]\right).
\]
However, this differs from the product of delta functions in Eq. \eqref{eq:9.6.5} only
by a factor $\operatorname{Det}\partial_3^2$, which although infinite is
field-independent and hence cancels in ratios like \eqref{eq:9.4.1}.} in Eq. \eqref{eq:9.6.5} to
enforce the constraints \eqref{eq:9.6.2} and \eqref{eq:9.6.3}.

The argument of the exponential in Eq. \eqref{eq:9.6.5} is evidently quadratic in the
independent components of $\boldsymbol\pi$ (say, $\pi_1$ and $\pi_2$), with
field-independent coefficients in the term of second order in
$\boldsymbol\pi$. Thus, according to Eq. \eqref{eq:9.A.9}, the integral over
$\boldsymbol\pi$ can be done (up to a constant factor) by setting
$\boldsymbol\pi$ equal to the stationary point of the argument of the
exponential, $\boldsymbol\pi=\dot{\mathbf a}$:
\begin{align}
\left\langle T\{\mathcal O_A\mathcal O_B\cdots\}\right\rangle_{\Omega}
={}&\int\left[\prod_{x,i}da_i(x)\prod_{x,\ell}d\psi_\ell(x)\right]
\notag\\
&\cdot\mathcal O_A\mathcal O_B\cdots
\notag\\
&\cdot\exp\Biggl\{i\int d^4x\Biggl[
 \frac12\dot{\mathbf a}^{\,2}
 -\frac12(\boldsymbol\nabla\times\mathbf a)^2
\notag\\
&\qquad+\mathbf a\cdot\mathbf j+\mathcal L_M\Biggr]
 -i\int dt\,V_{\mathrm{Coul}}+i\epsilon\text{ terms}\Biggr\}
\notag\\
&\cdot\left[\prod_x\delta(\boldsymbol\nabla\cdot\mathbf a(x))\right].
\tag{9.6.6}\label{eq:9.6.6}
\end{align}
To bring out the essential covariance of this result, we use a trick.
Introduce a new variable of integration $a^0(x)$, and replace the Coulomb term
$-\int dt\,V_{\mathrm{Coul}}$ in the action with
\begin{equation}
\int d^4x\left[-a^0(x)j^0(x)+\frac12
 (\boldsymbol\nabla a^0(x))^2\right].
\tag{9.6.7}\label{eq:9.6.7}
\end{equation}
Since \eqref{eq:9.6.7} is quadratic in $a^0$, the integral over $a^0$ can be done (up
to a constant factor) by setting $a^0(x)$ equal to the stationary point of
\eqref{eq:9.6.7}, i.e., to the solution of
\[
-j^0(x)-\boldsymbol\nabla^2a^0(x)=0
\]
or in other words, to
\begin{equation}
a^0(\mathbf{x},t)=\int d^3y\,
 \frac{j^0(\mathbf{y},t)}{4\pi|\mathbf{x}-\mathbf{y}|}.
\tag{9.6.8}\label{eq:9.6.8}
\end{equation}
Using this in Eq. \eqref{eq:9.6.7} just gives the Coulomb action
$-\int dt\,V_{\mathrm{Coul}}$. Hence we can rewrite the argument of the
exponential in Eq. \eqref{eq:9.6.6} as
\begin{align*}
&\frac12\dot{\mathbf a}^{,2}
 -\frac12(\boldsymbol\nabla\times\mathbf a)^2
 +\mathbf a\cdot\mathbf j+\mathcal L_M-a^0j^0
 +\frac12(\boldsymbol\nabla a^0)^2
\\
&\qquad=-\frac14f_{\mu\nu}f^{\mu\nu}+a_\mu j^\mu
 +\mathcal L_M+\text{total derivatives}
\end{align*}
with $f_{\mu\nu}=\partial_\mu a_\nu-\partial_\nu a_\mu$, and integrate over
$a^0$ as well as over $\mathbf a$ and matter fields. That is, the path
integral \eqref{eq:9.6.6} is now
\begin{align}
\left\langle T\{\mathcal O_A\mathcal O_B\cdots\}\right\rangle_{\Omega}
\propto{}&\int\left[\prod_{x,\mu}da_\mu(x)\right]
 \left[\prod_{x,\ell}d\psi_\ell(x)\right]
 \mathcal O_A\mathcal O_B\cdots
\notag\\
&\cdot\exp(iI[a,\psi])
 \prod_x\delta(\boldsymbol\nabla\cdot\mathbf a(x)),
\tag{9.6.9}\label{eq:9.6.9}
\end{align}
where $I$ is the original action
\begin{equation}
I[a,\psi]=\int d^4x\left[-\frac14f_{\mu\nu}f^{\mu\nu}
 +a_\mu j^\mu+\mathcal L_M\right]+i\epsilon\text{ terms}.
\tag{9.6.10}\label{eq:9.6.10}
\end{equation}

Now everything is manifestly Lorentz- and gauge-invariant, except for the
final product of delta functions which enforce the Coulomb gauge
condition.\footnote{Note that now $a^0(x)$ is \emph{not} equal to the value
\eqref{eq:9.6.8}, but is an independent variable of integration. We will not integrate
over $a^0(x)$ first, which would lead back to Eq. \eqref{eq:9.6.6}, but instead will
treat it in tandem with $\mathbf a(x)$.} To make further progress, we shall use
a simple version of a trick~[\hyperref[ref:9.4]{4},\hyperref[ref:9.5]{5}] that in Volume II will be used to treat the
more difficult case of non-Abelian gauge theories. For simplicity, we shall
deal here with the case where the operators $\mathcal O_A[A,\Psi]$,
$\mathcal O_B[A,\Psi],\ldots$ as well as the action $I[a,\psi]$ and measure
$[\prod da][\prod d\psi]$ are gauge-invariant.

First, replace the field variables of integration $a_\mu(x)$ and $\psi_\ell(x)$
everywhere in Eq. \eqref{eq:9.6.9} with the new variables
\begin{equation}
a_{\mu\Lambda}(x)\equiv a_\mu(x)+\partial_\mu\Lambda(x),
\tag{9.6.11}\label{eq:9.6.11}
\end{equation}
\begin{equation}
\psi_{\ell\Lambda}(x)\equiv
 \exp(iq_\ell\Lambda(x))\psi_\ell(x)
\tag{9.6.12}\label{eq:9.6.12}
\end{equation}
with arbitrary finite $\Lambda(x)$. This step is a mathematical triviality,
like changing an integral $\int_{-\infty}^{\infty}f(x)dx$ to read
$\int_{-\infty}^{\infty}f(y)dy$, and does not require use of the postulated
gauge invariance of the theory. Next, use gauge invariance to replace
$a_{\mu\Lambda}(x)$ and $\psi_{\ell\Lambda}(x)$ in the action, measure, and
$\mathcal O$-functions with the original fields $a_\mu(x)$ and $\psi_\ell(x)$,
respectively. Eq. \eqref{eq:9.6.9} then becomes
\begin{align}
&\left\langle T\{\mathcal O_A[A,\Psi],\mathcal O_B[A,\Psi],\ldots\}
 \right\rangle_{\Omega}
\notag\\
&\quad\propto\int\left[\prod_{x,\mu}da_\mu(x)\right]
 \left[\prod_{x,\ell}d\psi_\ell(x)\right]
 \mathcal O_A[a,\psi]\mathcal O_B[a,\psi]\cdots
\notag\\
&\qquad\cdot\exp(iI[a,\psi])
 \prod_x\delta(\boldsymbol\nabla\cdot\mathbf a(x)
 +\boldsymbol\nabla^2\Lambda(x)).
\tag{9.6.13}\label{eq:9.6.13}
\end{align}
Now, the function $\Lambda(x)$ was chosen at random, so despite appearances the
right-hand side of Eq. \eqref{eq:9.6.13} cannot depend on this function. We shall exploit
this fact to put the path integral in a much more convenient form. Multiply
Eq. \eqref{eq:9.6.13} by the functional
\begin{equation}
B[\Lambda,a]=\exp\left(-\frac12i\alpha\int d^4x\,
 [\partial_0a^0-\boldsymbol\nabla^2\Lambda]^2\right)
\tag{9.6.14}\label{eq:9.6.14}
\end{equation}
(where $\alpha$ is an arbitrary constant), and integrate over $\Lambda(x)$.
By shifting the integration variable $\Lambda(x)$, and noting the actual
$\Lambda$-independence of \eqref{eq:9.6.13}, we see that the effect is simply to
multiply Eq. \eqref{eq:9.6.13} with the field-independent constant
\begin{equation}
\int\left[\prod_xd\Lambda(x)\right]
 \exp\left(-\frac12i\alpha\int d^4x\,
 (\boldsymbol\nabla^2\Lambda)^2\right).
\tag{9.6.15}\label{eq:9.6.15}
\end{equation}
This factor cancels out in the connected part of the vacuum expectation value,
and thus has no physical effect. But \eqref{eq:9.6.13} is only $\Lambda$-independent
\emph{after} we integrate over $a^\mu(x)$ and $\psi(x)$. We can just as well
integrate over $\Lambda(x)$ \emph{before} we integrate over $a^\mu(x)$ and
$\psi(x)$, in which case the factor
$\prod_x\delta(\boldsymbol\nabla\cdot\mathbf a(x)
+\boldsymbol\nabla^2\Lambda)$ in Eq. \eqref{eq:9.6.13} is replaced with
\begin{align}
&\int\left[\prod_xd\Lambda(x)\right]
 \exp\left(-\frac12i\alpha\int d^4x\,
 [\partial_0a^0-\boldsymbol\nabla^2\Lambda]^2\right)
 \prod_x\delta(\boldsymbol\nabla\cdot\mathbf a(x)
 +\boldsymbol\nabla^2\Lambda)
\notag\\
&\qquad\propto\exp\left(-\frac12i\alpha\int d^4x\,
 (\partial_\mu a^\mu)^2\right),
\tag{9.6.16}\label{eq:9.6.16}
\end{align}
where ``$\propto$'' again means proportional with a field-independent factor.
Dropping constant factors, Eq. \eqref{eq:9.6.9} now becomes
\begin{align}
\left\langle T\{\mathcal O_A\mathcal O_B\cdots\}\right\rangle_{\Omega}
\propto{}&\int\left[\prod_{x,\mu}da_\mu(x)\right]
 \left[\prod_{x,\ell}d\psi_\ell(x)\right]
 \mathcal O_A\mathcal O_B\cdots\exp(iI_{\mathrm{eff}}[a,\psi]),
\tag{9.6.17}\label{eq:9.6.17}
\end{align}
where
\begin{equation}
I_{\mathrm{eff}}[a,\psi]=I[a,\psi]
 -\frac12\alpha\int(\partial_\mu a^\mu)^2d^4x.
\tag{9.6.18}\label{eq:9.6.18}
\end{equation}
This is now manifestly Lorentz-invariant.

We consider the new term in \eqref{eq:9.6.18} as a contribution to the unperturbed part
of the action, whose photonic part now reads
\begin{align}
I_0[a]
&=\int d^4x\left[-\frac14(\partial_\mu a_\nu-\partial_\nu a_\mu)
 (\partial^\mu a^\nu-\partial^\nu a^\mu)
 -\frac12\alpha(\partial_\mu a^\mu)^2+i\epsilon\text{ terms}\right]
\notag\\
&=-\frac12\int d^4x\,d^4y\,a^\mu(x)a^\nu(y)
 \mathcal D_{\mu x,\nu y},
\tag{9.6.19}\label{eq:9.6.19}
\end{align}
where
\begin{align}
\mathcal D_{\mu x,\nu y}
&=\left[\eta_{\mu\nu}
 \frac{\partial^2}{\partial x^\rho\partial y_\rho}
 -(1-\alpha)\frac{\partial^2}{\partial x^\mu\partial y^\nu}\right]
 \delta^4(x-y)+i\epsilon\text{ terms}
\notag\\
&=(2\pi)^{-4}\int d^4q\left[
 \eta_{\mu\nu}q^2-(1-\alpha)q_\mu q_\nu-i\epsilon\eta_{\mu\nu}\right]
 e^{iq\cdot(x-y)}.
\tag{9.6.20}\label{eq:9.6.20}
\end{align}
The photon propagator is then found immediately by inverting the $4\times4$
matrix in the integrand of Eq. \eqref{eq:9.6.20}
\begin{equation}
\Delta_{\mu x,\nu y}=(2\pi)^{-4}\int d^4q\left[
 \frac{\eta^{\mu\nu}}{q^2-i\epsilon}
 +\frac{1-\alpha}{\alpha}\frac{q^\mu q^\nu}{(q^2-i\epsilon)^2}\right]
 e^{iq\cdot(x-y)}.
\tag{9.6.21}\label{eq:9.6.21}
\end{equation}
We are free to choose $\alpha$ as seems most convenient. Two common choices
are $\alpha=1$, which yields the propagator in \emph{Feynman gauge}:
\begin{equation}
\Delta^{\mathrm{Feynman}}_{\mu x,\nu y}
=(2\pi)^{-4}\int d^4q\left[\frac{\eta^{\mu\nu}}{q^2-i\epsilon}\right]
 e^{iq\cdot(x-y)},
\tag{9.6.22}\label{eq:9.6.22}
\end{equation}
or $\alpha=\infty$, in which case the factor \eqref{eq:9.6.14} acts as a delta function,
and we obtain the propagator in \emph{Landau gauge} (often also called
\emph{Lorentz gauge}):
\begin{equation}
\Delta^{\mathrm{Landau}}_{\mu x,\nu y}
=(2\pi)^{-4}\int d^4q\left[
 \frac{\eta^{\mu\nu}}{q^2-i\epsilon}
 -\frac{q^\mu q^\nu}{(q^2-i\epsilon)^2}\right]e^{iq\cdot(x-y)}.
\tag{9.6.23}\label{eq:9.6.23}
\end{equation}
Practical calculations are made far more convenient by working with such
manifestly Lorentz-invariant interactions and propagators.
